\chapter{Ход работы}
\label{ch:chap2}

\section{Введение}

В предыдущей лабораторной работе я построил в MATLAB Simulink систему массового обслуживания M/M/1. СМО имела один генератор 
заявок, один сервер и одну очередь. В этой работе увеличим количество серверов до трех, а количество очередей оставим такое же.

\section{Модификация предыдущей схемы}

Модифицируем СМО из прошлой лабораторной работы следующим образом:

\begin{enumerate}
    \item Добавить 2 дополнительных сервера.
    \item Поставить на выходе из системы мультиплексор, который бы мультиплексировал выходы трех серверов.
    \item Добавить на выходе очереди демультиплексор, который будет "раскидывать" заявки по серверам.
    \item На входе очереди поставить мультиплексор на 4 входа, который бы мультиплексировал 3 обратных связи с серверов и выход генератора заявок.
\end{enumerate}

\section{Результат}

В результате модификации получена следующая схема:

\begin{figure}[H]
    \centering
    \includegraphics[width=1.0\textwidth]{smo.png}
    \caption{СМО с общей очередью}
\end{figure}

\section{Графики}

\begin{figure}[H]
    \centering
    \includegraphics[width=1.0\textwidth]{plot1.png}
    \caption{Графики}
\end{figure}

На графике в левом верхнем угле - график сгенерированных заявок. Остальные 3 графика - кол-во вышедших из серверов заявок (в том
числе отброшенные). Можно заметить, что нагрузка распределена равномерно (это связано с настройкой демультиплексора на выходе 
очереди, он работает по алгоритму round robin). Также можно обратить внимание, что сумма вышедших из серверов заявок больше, чем
число сгенерированных заявок. Это связано с тем, что сервера отбрасывают заявки с вероятностью 0.5.

\begin{figure}[H]
    \centering
    \includegraphics[width=1.0\textwidth]{plot2.png}
    \caption{График}
\end{figure}

На графике выше видно, что очередь выпустила почти 500 заявок.


\endinput
